\documentclass[12pt]{article}

\usepackage{amsmath}
\usepackage{amssymb}
\usepackage{amsfonts}
\usepackage{geometry}
\usepackage{setspace}
\usepackage[T1]{fontenc}
\usepackage[utf8]{inputenc}

\geometry{margin=1in}
\setstretch{1.2}

\begin{document}

\subsection*{Merton (1976): Option Pricing When Underlying Stock Returns Are Discontinuous}

\subsubsection*{Introduction}

The Black--Scholes model prices European options assuming stock prices follow a geometric
Brownian motion with continuous sample paths --- no sudden jumps. While Merton (1973) showed
the formula is robust to several other assumptions, continuity of the price path is essential:
it is what makes the perfect hedge possible. In practice, however, stock prices respond
discontinuously to earnings surprises, regulatory rulings, and macroeconomic shocks, producing
return distributions with far heavier tails than the log-normal model predicts. Merton (1976)
extends the framework by incorporating a jump component directly into the price process.

\subsubsection*{The Jump-Diffusion Model}

The stock price dynamics are governed by:

\begin{equation}
    \frac{dS_t}{S_t} = (\alpha - \lambda k)\,dt + \sigma\,dZ_t + dq_t
\end{equation}

where $\alpha$ is the expected return, $\sigma^2$ the diffusion variance (conditional on no
jump), $dZ_t$ a standard Wiener process, and $dq_t$ a Poisson jump component with intensity
$\lambda$ (mean jumps per unit time). On each jump the stock moves from $S_t$ to $S_t Y$,
where $Y>0$ is a random multiplier with $k \equiv \mathbb{E}(Y-1)$; the $\lambda k$ correction
keeps the total expected return equal to $\alpha$. Setting $\lambda = 0$ recovers Black--Scholes
exactly. The Poisson process is the natural choice for jump timing: its memoryless property
reflects the informational efficiency requirement that the arrival of market-moving news be
unpredictable.

\subsubsection*{Option Pricing and the Breakdown of the Black--Scholes Hedge}

When a jump occurs, the stock price moves discontinuously and the option value responds
nonlinearly. Because portfolio replication is a linear operation, no combination of stock and
bond can hedge this nonlinear jump exposure, making the market \textit{incomplete}. Merton
resolves the indeterminacy by assuming jump risk is firm-specific and therefore
non-systematic (diversifiable). Under the CAPM, a zero-beta portfolio must earn the risk-free
rate, which yields the integro-differential pricing equation:

\begin{equation}
    0 = \tfrac{1}{2}\sigma^2 S^2 F_{SS} + (r - \lambda k)\,S F_S - F_\tau - rF
      + \lambda\,\mathbb{E}\bigl\{F(SY,\tau) - F(S,\tau)\bigr\}
\end{equation}

with $F(0,\tau)=0$ and $F(S,0)=\max[0,\,S-E]$. The equation is preference-free and reduces to
the Black--Scholes PDE when $\lambda=0$. For log-normally distributed jump sizes with
$\mathbb{E}(Y)=1$ and $\mathrm{Var}(\ln Y)=\delta^2$, the closed-form solution is:

\begin{equation}
    F(S,\tau) = \sum_{n=0}^{\infty}
    \frac{e^{-\lambda\tau}(\lambda\tau)^n}{n!}\;
    W\!\left(S,\tau;\,E,\,\sigma^2+\frac{n\delta^2}{\tau},\,r\right)
\end{equation}

where $W(\cdot)$ is the Black--Scholes formula. The price is a Poisson-weighted average of
Black--Scholes prices, one for each possible number of jumps $n$, each evaluated at an inflated
variance $\sigma^2 + n\delta^2/\tau$ that reflects the additional uncertainty from $n$ jumps.

\subsubsection*{Implications, Strengths, and Limitations}

Equation~(3) explains the \textit{volatility smile}: deep out-of-the-money options are
underpriced by Black--Scholes because a single large jump can push the stock past the strike
instantly, a scenario a continuous diffusion nearly cannot produce over short horizons. This
generates the well-known pattern of implied volatilities that rise away from the money. More
broadly, the jump-diffusion process produces excess kurtosis consistent with empirical return
distributions.

The model's key strengths are its analytical tractability (the series converges rapidly and
each term is a Black--Scholes price), its clear economic interpretation (Poisson arrivals
reflect unpredictable firm-specific news), and its role as the foundation for a large family
of later models. Its main limitations are the constant jump intensity $\lambda$ (unable to
capture time-varying crash risk or volatility clustering), the symmetric log-normal jump
distribution (unable to match the steep downside skew of index options), and the absence of
stochastic volatility, all of which are addressed in extensions such as Bates (1996).

\subsubsection*{Relevance to QQQ Options}

The QQQ ETF tracks the NASDAQ-100, a technology-heavy index particularly susceptible to
jump-like behaviour. The implied volatility surface extracted from QQQ options displays a
pronounced downside skew --- out-of-the-money puts carry significantly higher implied
volatilities than at-the-money options --- consistent with equation~(3). We calibrate
$(\sigma,\lambda,\delta^2)$ to the observed QQQ cross-section and compare the fit to
Black--Scholes. The jump-diffusion model is expected to outperform for short-maturity,
far-from-the-money options where jump risk is most material, and to converge to the
Black--Scholes fit for longer-dated near-the-money contracts.

\end{document}